\documentclass[10pt]{article}
\usepackage[T1]{fontenc}
\usepackage[utf8]{inputenc}
\usepackage[french]{babel}
\usepackage{amsmath,amssymb,mathrsfs,mathpazo}
\usepackage{geometry}
\geometry{a4paper,margin=2cm}
\usepackage{graphicx,array,tabularx,tabularray,colortbl,multirow,multicol,booktabs,enumitem,pifont,ulem,fancybox,diagbox,cancel}
\usepackage{pstricks,pstricks-add,pst-plot,pst-tree,pst-node,pst-text,pst-func,pst-math,pst-bspline,pst-3dplot,pst-eucl,pst-solides3d}
\let\pscalculate\undefined
\usepackage{pst-bezier}
\usepackage{tikz,pgfplots,tkz-tab}
\usetikzlibrary{arrows,arrows.meta,patterns,calc,positioning,trees,decorations.markings,shapes.geometric,shapes.misc,shapes.arrows,intersections,decorations.pathmorphing,decorations.pathreplacing,backgrounds,shadows,babel}
\pgfplotsset{compat=1.18}
\usepgfplotslibrary{fillbetween}
\usepackage[np]{numprint}
\usepackage{listings,makecell,pas-tableur,twoopt}
\let\barre\undefined
\usepackage[pstricks]{tablvar}
\definecolor{darkgreen}{RGB}{0,100,0}
\providecommand{\red}{\color{red}}
\providecommand{\blue}{\color{blue}}
\providecommand{\green}{\color{green}}
\providecommand{\black}{\color{black}}
\providecommand{\white}{\color{white}}
\providecommand{\gray}{\color{gray}}
\providecommand{\cyan}{\color{cyan}}
\providecommand{\magenta}{\color{magenta}}
\providecommand{\yellow}{\color{yellow}}
\providecommand{\violet}{\color{violet}}
\providecommand{\orange}{\color{orange}}
\providecommand{\euro}{\texteuro}
\providecommand{\EUR}{\texteuro}
\providecommand{\ieme}{\textsuperscript{e}}
\providecommand{\degres}{\ensuremath{^\circ}}
\providecommand{\degree}{\ensuremath{^\circ}}
\providecommand{\vv}[1]{\overrightarrow{#1}}
\providecommand{\Coord}[1]{#1}
\providecommand{\decosix}{}
\providecommand{\starredbullet}{\ensuremath{\bullet}}
\pagestyle{empty}
\setlength{\parindent}{0pt}
\usepackage{adjustbox}\DeclareUnicodeCharacter{2236}{:}\DeclareUnicodeCharacter{2190}{\ensuremath{\leftarrow}}
\providecommand{\R}{\mathbb{R}}
\providecommand{\N}{\mathbb{N}}
\providecommand{\Z}{\mathbb{Z}}
\providecommand{\Q}{\mathbb{Q}}
\providecommand{\D}{\mathbb{D}}
\providecommand{\C}{\mathbb{C}}
\providecommand{\e}{{\mathrm{e}}}
\providecommand{\vect}[1]{\overrightarrow{#1}}
\providecommand{\vectt}[1]{\overrightarrow{\mathrm{#1}}}
\providecommand{\vvt}[1]{\overrightarrow{\mathrm{#1}}}
\providecommand{\np}[2][]{{#2~#1}}
\providecommand{\vv}[1]{{\overrightarrow{#1}}}
\renewcommand{\d}{{\mathrm{d}}}
\providecommand{\degres}{{^\circ}}
\providecommand{\degre}{{^\circ}}
\providecommand{\blue}{\color{blue}}
\providecommand{\red}{\color{red}}
\providecommand{\enskip}{\quad}
\providecommand{\llbracket}{\lbrack\!\lbrack}
\providecommand{\rrbracket}{\rbrack\!\rbrack}
\providecommand{\euro}{{\text{€}}}
\providecommand{\no}{{\text{n°}}}
\providecommand{\interval}[2]{{[#1;#2]}}
\providecommand{\intervFO}[2]{{[#1;#2[}}
\providecommand{\dotfill}{\ldots}

\providecommand{\ds}{\displaystyle}
\providecommand{\textding}[1]{#1}
\providecommand{\Oij}{\(\mathrm{O};\vec{\imath},\vec{\jmath}\)}
\providecommand{\Oijk}{\(\mathrm{O};\vec{\imath},\vec{\jmath},\vec{k}\)}
\providecommand{\Ouv}{\(\mathrm{O};\vec{u},\vec{v}\)}
\providecommand{\intervFF}[2]{\ensuremath{[#1;#2]}}
\renewcommand{\euro}{\ifmmode\text{\texteuro}\else\texteuro\fi}
\DeclareUnicodeCharacter{00B0}{\ifmmode^{\circ}\else\textdegree\fi}
\DeclareUnicodeCharacter{20AC}{\ifmmode\text{\texteuro}\else\texteuro\fi}
% Énoncés : Éducation nationale. Transcription et corrigés : APMEP, Denis Vergès et contributeurs. Découpage et classement : CoopMaths, Sébastien Lozano, Éric Elter et contributeurs. Adaptation de la présentation : Maths Explorer.
% CC BY-SA : https://coopmaths.fr/www/about/
% https://coopmaths.fr/alea/static/eam/2025/tex/eam_specifique_2025_06_sujet0v2_2.tex
\begin{document}
\noindent\textit{APMEP / CoopMaths. CC BY-SA. Présentation adaptée par Maths Explorer.}\par\medskip
On étudie la croissance d'une population de champignons.

\subsubsection*{Partie A}
Au début de l'expérience, on dispose de 100 champignons. Toutes les dix minutes, on mesure l'évolution de leur nombre. On obtient les résultats suivants :

\begin{minipage}[t]{0.5\textwidth}
\centering

\begin{tabular}{|c|c|}
\hline
\makecell{Temps écoulé \\ (en minutes)} & \makecell{Nombre de \\ champignons} \\
\hline
0 & 100 \\
10 & 125 \\
20 & 150 \\
30 & 175 \\
\hline
\end{tabular}

\end{minipage}
\hfill
\begin{minipage}[t]{0.45\textwidth}
\vspace{-1.5cm}
\centering

\begin{tikzpicture}[scale=0.1]
\foreach \x in {0,10,...,35} {
\draw[gray] (\x,0) -- (\x,40);
}
\foreach \y in {0,50,...,190} {
\draw[gray] (0,\y/5) -- (35,\y/5);
}
\foreach \x in {0,2,...,35} {
\draw[lightgray] (\x,0) -- (\x,40);
}
\foreach \y in {0,10,...,190} {
\draw[lightgray] (0,\y/5) -- (35,\y/5);
}
\draw[->, line width=1pt] (0,0) -- (35,0) node[]{};
\draw[->, line width=1pt] (0,0) -- (0,40) node[]{};
\node[align=right, font=\tiny] at (28,5) {temps écoulé\\ (en minutes)};
\node[align=left, font=\tiny] at (8,37) {nombre de\\ champignons};
\foreach \x in {0,10,20,30} {
\draw (\x,0) node[below]{$\x$};
}
\foreach \y in {50,100,150} {
\draw (0,\y/5) node[left]{$\y$};
}
\foreach \x/\y in {0/100, 10/125, 20/150, 30/175} {
\filldraw[red] (\x,\y/5) circle (15pt);
}
\end{tikzpicture}

\end{minipage}

Soit $n$ un entier naturel. On note $u_n$ le nombre de champignons après $n$ périodes de \textbf{dix} minutes. Ainsi $u_0 = 100$, $u_1 = 125$, $u_2 = 150$...

\begin{enumerate}
\item Justifier que les termes $u_0$, $u_1$, $u_2$, $u_3$ sont en progression arithmétique.

\item En supposant que la population de champignons continue d'évoluer selon le même rythme, montrer qu'elle aura quadruplé deux heures après le début de l'expérience.
\end{enumerate}

\begin{minipage}[t]{0.6\textwidth}
\subsubsection*{Partie B}
En réalité, on constate que la population de champignons a quadruplé 80 minutes après le début de l'expérience. De nouvelles mesures donnent les résultats suivants.
\end{minipage}
\hfill
\begin{minipage}[t]{0.35\textwidth}
\vspace{0pt}
\begin{center}
\begin{tabular}{|c|c|}
\hline
\makecell{Temps écoulé \\ (en minutes)} & \makecell{Nombre de \\ champignons} \\
\hline
0 & 100 \\
40 & 200 \\
80 & 400 \\
120 & 800 \\
\hline
\end{tabular}
\end{center}
\end{minipage}

Soit $n$ un entier naturel. On note $v_n$ le nombre de champignons, après $n$ périodes de \textbf{quarante} minutes. Ainsi $v_0 = 100$, $v_1 = 200$, $v_2 = 400$...

\begin{enumerate}
\item Montrer que les termes $v_0$, $v_1$, $v_2$, $v_3$ sont en progression géométrique.

\item On suppose que la suite $(v_n)$ est une suite géométrique de raison 2.

Indiquer sans justifier lequel des 4 graphiques ci-dessous est susceptible de représenter la suite $(v_n)$.
\end{enumerate}

\begin{minipage}[t]{0.24\textwidth}
\begin{center}
Graphique 1 \\
\begin{tikzpicture}
\draw[semithick,->,>=latex] (0,0) -- (3.5,0) ;
\draw[semithick,->,>=latex] (0,0) -- (0,2) ;
\foreach \coord in {(0.5,0.1),(1,0.2),(1.5,0.3),(2,0.5),(2.5,1),(3,1.9)} {
        \fill[red] \coord circle (2pt);
    }
\end{tikzpicture}
\end{center}
\end{minipage}
\hfill
\begin{minipage}[t]{0.24\textwidth}
\begin{center}
Graphique 2 \\
\begin{tikzpicture}
\draw[semithick,->,>=latex] (0,0) -- (3.5,0) ;
\draw[semithick,->,>=latex] (0,0) -- (0,2) ;
\foreach \coord in {(0.1,0.3),(0.2,0.5),(0.45,0.8),(0.9,1),(1.8,1.3),(3.5,1.6)} {
        \fill[red] \coord circle (2pt);
    }
\end{tikzpicture}
\end{center}
\end{minipage}
\hfill
\begin{minipage}[t]{0.24\textwidth}
\begin{center}
Graphique 3 \\
\begin{tikzpicture}
\draw[semithick,->,>=latex] (0,0) -- (3.5,0) ;
\draw[semithick,->,>=latex] (0,0) -- (0,2) ;
\foreach \coord in {(0.5,0.25),(1,0.5),(1.5,0.75),(2,1),(2.5,1.25),(3,1.5)} {
        \fill[red] \coord circle (2pt);
    }
\end{tikzpicture}
\end{center}
\end{minipage}
\hfill
\begin{minipage}[t]{0.24\textwidth}
\begin{center}
Graphique 4 \\
\begin{tikzpicture}
\draw[semithick,->,>=latex] (0,0) -- (3.5,0) ;
\draw[semithick,->,>=latex] (0,0) -- (0,2) ;
\foreach \coord in {(0.5,1.9),(1,1),(1.5,0.65),(2,0.45),(2.5,0.4),(3,0.3)} {
        \fill[red] \coord circle (2pt);
    }
\end{tikzpicture}
\end{center}
\end{minipage}

\begin{center}
\begin{minipage}{0.75\textwidth}
\begin{enumerate}[start=3]
\item Quel sera le nombre de champignons quatre heures après le début de l'expérience ?

\item Cinq heures après le début de l'expérience, on dénombre environ 18 000 champignons. Est-ce cohérent avec le modèle choisi ?
\end{enumerate}
\end{minipage}
\hfill
\fbox{
\begin{minipage}{0.2\linewidth}
\centering
Aide au calcul :

$2^6 = 64$

$2^7 = 128$

$2^8 = 256$

$2^9 = 512$

$2^{10} = 1024$
\end{minipage}
}
\end{center}

\end{document}
